% ---------------------------------------------------------------------------
% teste-visual.tex — prova visual da classe doeapostila.
% Exercita todos os elementos previstos para os cinco cadernos: hierarquia de
% títulos, equações numeradas com as macros de notação, tabela booktabs com
% siunitx em vírgula decimal, figuras TikZ e PGFPlots com legenda abaixo,
% o ambiente aplicacao, listas do enumitem, acentuação portuguesa plena e a
% página de direitos autorais.
%
% Compilação (MiKTeX fora do PATH):
%   $env:PATH = "C:\Users\Narcísius\AppData\Local\Programs\MiKTeX\miktex\bin\x64;" + $env:PATH
%   pdflatex -interaction=nonstopmode teste-visual.tex   (duas vezes)
% ---------------------------------------------------------------------------
\documentclass{doeapostila}

% ---------------------------------------------------------------------------
% macros.tex — notação matemática comum aos cinco cadernos.
% Fixada antes de qualquer caderno começar, para os cinco não divergirem.
% Uso: % ---------------------------------------------------------------------------
% macros.tex — notação matemática comum aos cinco cadernos.
% Fixada antes de qualquer caderno começar, para os cinco não divergirem.
% Uso: % ---------------------------------------------------------------------------
% macros.tex — notação matemática comum aos cinco cadernos.
% Fixada antes de qualquer caderno começar, para os cinco não divergirem.
% Uso: \input{../estilo/macros} no preâmbulo de cada main.tex.
% ---------------------------------------------------------------------------

% Somas de quadrados
\newcommand{\SQT}{\mathrm{SQ}_{\mathrm{T}}}    % total
\newcommand{\SQR}{\mathrm{SQ}_{\mathrm{R}}}    % da regressão
\newcommand{\SQr}{\mathrm{SQ}_{\mathrm{r}}}    % residual
\newcommand{\SQfaj}{\mathrm{SQ}_{\mathrm{faj}}}% falta de ajuste
\newcommand{\SQep}{\mathrm{SQ}_{\mathrm{ep}}}  % erro puro

% Médias quadráticas: \MQ{R}, \MQ{r}, \MQ{faj}, \MQ{ep}
\newcommand{\MQ}[1]{\mathrm{MQ}_{\mathrm{#1}}}

% Coeficientes de determinação
\newcommand{\Rdois}{R^{2}}
\newcommand{\Rajus}{R^{2}_{\mathrm{aj}}}

% Efeitos: \efeito{1}, \efeito{12}
\newcommand{\efeito}[1]{E_{#1}}

% Álgebra matricial do ajuste por mínimos quadrados
\newcommand{\vetbeta}{\boldsymbol{\beta}}
\newcommand{\matX}{\mathbf{X}}
\newcommand{\vety}{\mathbf{y}}
 no preâmbulo de cada main.tex.
% ---------------------------------------------------------------------------

% Somas de quadrados
\newcommand{\SQT}{\mathrm{SQ}_{\mathrm{T}}}    % total
\newcommand{\SQR}{\mathrm{SQ}_{\mathrm{R}}}    % da regressão
\newcommand{\SQr}{\mathrm{SQ}_{\mathrm{r}}}    % residual
\newcommand{\SQfaj}{\mathrm{SQ}_{\mathrm{faj}}}% falta de ajuste
\newcommand{\SQep}{\mathrm{SQ}_{\mathrm{ep}}}  % erro puro

% Médias quadráticas: \MQ{R}, \MQ{r}, \MQ{faj}, \MQ{ep}
\newcommand{\MQ}[1]{\mathrm{MQ}_{\mathrm{#1}}}

% Coeficientes de determinação
\newcommand{\Rdois}{R^{2}}
\newcommand{\Rajus}{R^{2}_{\mathrm{aj}}}

% Efeitos: \efeito{1}, \efeito{12}
\newcommand{\efeito}[1]{E_{#1}}

% Álgebra matricial do ajuste por mínimos quadrados
\newcommand{\vetbeta}{\boldsymbol{\beta}}
\newcommand{\matX}{\mathbf{X}}
\newcommand{\vety}{\mathbf{y}}
 no preâmbulo de cada main.tex.
% ---------------------------------------------------------------------------

% Somas de quadrados
\newcommand{\SQT}{\mathrm{SQ}_{\mathrm{T}}}    % total
\newcommand{\SQR}{\mathrm{SQ}_{\mathrm{R}}}    % da regressão
\newcommand{\SQr}{\mathrm{SQ}_{\mathrm{r}}}    % residual
\newcommand{\SQfaj}{\mathrm{SQ}_{\mathrm{faj}}}% falta de ajuste
\newcommand{\SQep}{\mathrm{SQ}_{\mathrm{ep}}}  % erro puro

% Médias quadráticas: \MQ{R}, \MQ{r}, \MQ{faj}, \MQ{ep}
\newcommand{\MQ}[1]{\mathrm{MQ}_{\mathrm{#1}}}

% Coeficientes de determinação
\newcommand{\Rdois}{R^{2}}
\newcommand{\Rajus}{R^{2}_{\mathrm{aj}}}

% Efeitos: \efeito{1}, \efeito{12}
\newcommand{\efeito}[1]{E_{#1}}

% Álgebra matricial do ajuste por mínimos quadrados
\newcommand{\vetbeta}{\boldsymbol{\beta}}
\newcommand{\matX}{\mathbf{X}}
\newcommand{\vety}{\mathbf{y}}


\definecaderno{0}{Prova visual da classe}

\begin{document}

\begin{center}
  {\large\bfseries Prova visual da classe \texttt{doeapostila}}\par\smallskip
  {\normalsize Murilo Miguel Narciso}\par\smallskip
  {\footnotesize Instituto Militar de Engenharia \textbullet{} \today}
\end{center}

\vspace{0.6\baselineskip}

\noindent
Este documento não tem valor didático. Ele existe para que cada recurso da
classe seja visto uma vez antes de os cadernos começarem a ser escritos, e
para que qualquer regressão tipográfica introduzida depois seja detectada numa
única compilação.

\vspace{0.6\baselineskip}

\tableofcontents

\section{Fundamentos da variação}

A afirmação de que uma alteração no processo produziu um ganho só se sustenta
quando a variação observada excede aquela que o próprio procedimento
experimental gera sem que nada seja mudado. É essa comparação, e não a
magnitude absoluta da diferença, que fundamenta a inferência. A avaliação
depende, portanto, de uma estimativa independente do erro experimental, obtida
por réplicas genuínas e não por medições repetidas de uma mesma amostra.

\subsection{Decomposição da soma de quadrados}

A variabilidade total da resposta é separada em uma parcela explicada pelo
modelo e outra atribuída ao resíduo, conforme a
Equação~\eqref{eq:decomposicao}:

\begin{equation}
  \SQT = \SQR + \SQr .
  \label{eq:decomposicao}
\end{equation}

O resíduo, por sua vez, admite subdivisão sempre que houver réplicas nos
mesmos pontos experimentais. A parcela de erro puro mede a dispersão interna
aos pontos replicados e independe do modelo adotado; a parcela de falta de
ajuste mede o quanto o modelo se afasta das médias observadas:

\begin{equation}
  \SQr = \SQfaj + \SQep ,
  \qquad
  F_{\mathrm{faj}} = \frac{\MQ{faj}}{\MQ{ep}} .
  \label{eq:faltadeajuste}
\end{equation}

\subsection{Coeficientes de determinação}

O coeficiente de determinação e sua versão ajustada aos graus de liberdade são
definidos na Equação~\eqref{eq:erredois}. A distinção é relevante porque o
primeiro nunca decresce com a inclusão de termos, ainda que supérfluos, ao
passo que o segundo penaliza a perda de graus de liberdade:

\begin{equation}
  \Rdois = 1 - \frac{\SQr}{\SQT} ,
  \qquad
  \Rajus = 1 - \frac{\MQ{r}}{\MQ{T}} .
  \label{eq:erredois}
\end{equation}

\subsubsection{Estimativa dos coeficientes}

Os coeficientes do modelo são obtidos por mínimos quadrados. Em notação
matricial, com $\matX$ a matriz de planejamento e $\vety$ o vetor de respostas,
a estimativa é dada pela Equação~\eqref{eq:minimosquadrados}:

\begin{equation}
  \hat{\vetbeta} = \left( \matX^{\top} \matX \right)^{-1} \matX^{\top} \vety .
  \label{eq:minimosquadrados}
\end{equation}

\subsubsection{Efeito de um fator}

No planejamento fatorial de dois níveis, o efeito principal do fator $i$ é a
diferença entre a média das respostas no nível superior e a média no nível
inferior, de modo que $\efeito{i}$ corresponde ao dobro do coeficiente
homônimo do modelo de regressão nas variáveis codificadas. A interação de
segunda ordem $\efeito{12}$ segue a mesma construção, com os sinais obtidos
pelo produto das colunas correspondentes.

\section{Apresentação de resultados}

\subsection{Tabela de efeitos}

A Tabela~\ref{tab:efeitos} reúne os efeitos estimados em um planejamento
$2^{3}$ conduzido com duas réplicas. A vírgula é o separador decimal, conforme
a norma brasileira, e o alinhamento pela casa decimal é obrigatório em todas as
tabelas numéricas dos cadernos.

\begin{table}[H]
  \centering
  \caption{Efeitos estimados no planejamento $2^{3}$ com duas réplicas.
    Erro padrão comum de \num{1.80} unidades, calculado a partir da variância
    combinada das réplicas.}
  \label{tab:efeitos}
  \begin{tabular}{l S[table-format=-2.2] S[table-format=1.2] S[table-format=-2.2]}
    \toprule
    {Termo} & {Efeito} & {Erro padrão} & {Razão $t$} \\
    \midrule
    Média                          & 67.30  & 0.90 & 74.78 \\
    Temperatura                    & 22.90  & 1.80 & 12.72 \\
    Catalisador                    & -13.90 & 1.80 & -7.72 \\
    Concentração                   & 8.90   & 1.80 & 4.94  \\
    Temperatura $\times$ catalisador & -8.60  & 1.80 & -4.78 \\
    \bottomrule
  \end{tabular}
\end{table}

Números fora de tabela também seguem a convenção: a variância combinada vale
\num{6.5}, o rendimento acumulado atingiu \num{12345.678} unidades e a
resistência característica ficou em \qty{2.95}{\mega\pascal} aos vinte e oito
dias.

\subsection{Representação geométrica do planejamento}

A Figura~\ref{fig:quadrado} representa o planejamento $2^{2}$ como um quadrado
no espaço das variáveis codificadas. Cada vértice corresponde a uma combinação
de níveis, e o contraste que define um efeito principal é a diferença entre as
médias de dois lados opostos.

\begin{figure}[H]
  \centering
  \begin{tikzpicture}[scale=2.6, line width=0.5pt]
    \draw (0,0) rectangle (1,1);
    \draw[->, line width=0.4pt] (-0.28,0) -- (1.28,0) node[right, font=\footnotesize] {$x_{1}$};
    \draw[->, line width=0.4pt] (0,-0.28) -- (0,1.28) node[above, font=\footnotesize] {$x_{2}$};
    \foreach \x/\y/\rot in {0/0/below left, 1/0/below right, 0/1/above left, 1/1/above right}
      \fill (\x,\y) circle (1.4pt);
    \node[below left, font=\footnotesize]  at (0,0) {$(-,-)$};
    \node[below right, font=\footnotesize] at (1,0) {$(+,-)$};
    \node[above left, font=\footnotesize]  at (0,1) {$(-,+)$};
    \node[above right, font=\footnotesize] at (1,1) {$(+,+)$};
    \fill (0.5,0.5) circle (1.4pt);
    \node[right, font=\footnotesize] at (0.54,0.5) {ponto central};
  \end{tikzpicture}
  \caption{Planejamento fatorial $2^{2}$ com ponto central, em variáveis
    codificadas. O ponto central não altera as estimativas dos efeitos
    principais e fornece o grau de liberdade necessário ao teste de curvatura.}
  \label{fig:quadrado}
\end{figure}

\subsection{Ajuste linear e ajuste quadrático}

A Figura~\ref{fig:ajustes} contrasta dois modelos ajustados ao mesmo conjunto
de observações. A cor de apoio é reservada à série ajustada; os dados
permanecem em preto.

\begin{figure}[H]
  \centering
  \begin{subfigure}{0.48\linewidth}
    \centering
    \begin{tikzpicture}
      \begin{axis}[
        width=\linewidth, height=4.6cm,
        xlabel={$x$}, ylabel={$y$},
        xmin=25, xmax=75, ymin=0, ymax=60,
        tick label style={font=\scriptsize},
        label style={font=\footnotesize},
        grid=major, grid style={gray!25, line width=0.3pt},
        axis line style={black, line width=0.4pt},
      ]
        \addplot[only marks, mark=*, mark size=1.3pt, black]
          coordinates {(30,17) (30,21) (40,32) (40,36) (50,43) (50,47)
                       (60,48) (60,52) (70,43) (70,47)};
        \addplot[azulDoE, line width=0.9pt, domain=28:72, samples=2]
          {13.6 + 0.51*x};
      \end{axis}
    \end{tikzpicture}
    \caption{Modelo linear, $\Rdois = \num{0.777873}$.}
    \label{fig:ajuste-linear}
  \end{subfigure}
  \hfill
  \begin{subfigure}{0.48\linewidth}
    \centering
    \begin{tikzpicture}
      \begin{axis}[
        width=\linewidth, height=4.6cm,
        xlabel={$x$}, ylabel={$y$},
        xmin=25, xmax=75, ymin=0, ymax=60,
        tick label style={font=\scriptsize},
        label style={font=\footnotesize},
        grid=major, grid style={gray!25, line width=0.3pt},
        axis line style={black, line width=0.4pt},
      ]
        \addplot[only marks, mark=*, mark size=1.3pt, black]
          coordinates {(30,17) (30,21) (40,32) (40,36) (50,43) (50,47)
                       (60,48) (60,52) (70,43) (70,47)};
        \addplot[azulDoE, line width=0.9pt, domain=28:72, samples=60]
          {-40.9 + 2.66*x - 0.0215*x^2};
      \end{axis}
    \end{tikzpicture}
    \caption{Modelo quadrático, $\Rdois = \num{0.993459}$.}
    \label{fig:ajuste-quadratico}
  \end{subfigure}
  \caption{Comparação entre o ajuste linear e o ajuste quadrático de um mesmo
    conjunto de dez observações. O erro puro é idêntico nos dois casos, porque
    depende apenas das réplicas; o que muda é a falta de ajuste.}
  \label{fig:ajustes}
\end{figure}

\begin{aplicacao}{Diluição de erros de laboratório}
  Em um ensaio de compressão simples conduzido em duplicata, a leitura de um
  dos corpos de prova foi comprometida por um posicionamento incorreto da
  célula de carga. Como o planejamento previa réplicas genuínas, o desvio
  ficou confinado à estimativa de erro puro e não contaminou os efeitos
  estimados. A razão $F_{\mathrm{faj}} = \MQ{faj} / \MQ{ep}$ caiu de
  \num{55.4} para \num{0.447} após a inclusão dos termos quadráticos, sem que
  o erro puro se alterasse, o que confirma que a deficiência residia no
  modelo e não nos dados.
\end{aplicacao}

\section{Procedimento recomendado}

\subsection{Sequência de execução}

A ordem abaixo é a que reduz o número de ensaios necessários quando o material
disponível é escasso:

\begin{enumerate}[label=\arabic*., itemsep=2pt, topsep=4pt, leftmargin=*]
  \item Definir a resposta de interesse e a precisão exigida na sua medição.
  \item Levantar os fatores candidatos e distinguir os controláveis dos
    ruídos, registrando as razões de cada classificação.
  \item Executar uma triagem fracionada, aceitando confundimento entre
    interações de ordem elevada.
  \item Acrescentar réplicas apenas nos pontos que sustentam a estimativa de
    erro puro.
  \item Ampliar para composição central somente se o teste de curvatura
    indicar que o modelo de primeira ordem é insuficiente.
\end{enumerate}

\subsection{Condições que invalidam a análise}

\begin{itemize}[label=\textbullet, itemsep=2pt, topsep=4pt, leftmargin=*]
  \item Ausência de aleatorização na ordem de execução dos ensaios.
  \item Repetição de medição tomada como réplica, o que subestima o erro
    experimental e infla artificialmente a significância dos efeitos.
  \item Extrapolação da superfície ajustada para além da região experimental.
\end{itemize}

\subsection{Vocabulário e acentuação}

O texto dos cadernos é escrito em português brasileiro com acentuação plena.
Esta seção existe para que a codificação seja verificada de fato:
ação, avaliação, execução, precisão, confundimento, coeficiente, variância,
influência, análise, ângulo, três, é, ê, á, à, ó, ú, í, â, ç, ã, õ.
Composição central rotacionável, réplicas genuínas, região experimental,
não-aditividade, mínimos quadrados, propagação de incerteza, hipótese nula.
Não há sinônimos em inglês no corpo do texto, e as unidades seguem o Sistema
Internacional.

A plataforma que acompanha estas apostilas está em
\url{https://research.iacommurilo.com.br}, e o índice de todo o material pode
ser consultado a partir da Seção~\ref{sec:encerramento}.

\section{Encerramento}
\label{sec:encerramento}

\begin{description}[style=nextline, leftmargin=1.4em, itemsep=3pt, topsep=4pt]
  \item[Erro puro] Estimativa da variabilidade obtida exclusivamente de
    réplicas genuínas nos mesmos pontos experimentais.
  \item[Falta de ajuste] Parcela do resíduo atribuída à inadequação da forma
    funcional do modelo, disponível apenas quando existe erro puro.
\end{description}

Encerrada a verificação, a página seguinte reproduz a nota de direitos
autorais que fecha cada um dos cinco cadernos.

\paginadireitos

\end{document}
